\documentclass[border=5pt]{standalone}
\usepackage[T1]{fontenc}
\usepackage{tikz}
\usepackage{luacode}

\definecolor{garnet}{HTML}{73000A}
\definecolor{atlantic}{HTML}{466A9F}
\definecolor{congaree}{HTML}{1F414D}
\definecolor{warmgrey}{HTML}{676156}
\definecolor{bk70}{HTML}{5C5C5C}
\definecolor{bk30}{HTML}{C7C7C7}

\begin{document}
\begin{tikzpicture}

\begin{luacode}
-- ════════════════════════════════════════════════════════════
-- Lens shape showcase: one number (R) per surface, auto thickness
-- ════════════════════════════════════════════════════════════
-- Sign convention:
--   R > 0 : surface bulges LEFT  (convex to incoming light)
--   R < 0 : surface bulges RIGHT (concave to incoming light)
--   |R| large (e.g. 999) : nearly flat

local min_wall = 0.08   -- minimum glass thickness anywhere
local h = 1.0           -- half-aperture for all examples

-- ── Auto-thickness calculation ──
local function auto_t(Rl, Rr, h)
    local absRl = math.abs(Rl)
    local absRr = math.abs(Rr)
    local sag_l = (absRl >= h) and (absRl - math.sqrt(absRl*absRl - h*h)) or absRl
    local sag_r = (absRr >= h) and (absRr - math.sqrt(absRr*absRr - h*h)) or absRr
    -- Convex surfaces make edges thinner; concave make edges thicker
    local dl = (Rl > 0) and -sag_l or sag_l
    local dr = (Rr < 0) and -sag_r or sag_r
    local edge_delta = dl + dr  -- negative = edges thinner than center
    -- t = center thickness; ensure min_wall everywhere
    local t = min_wall - math.min(0, edge_delta)
    if t < min_wall then t = min_wall end
    return t
end

-- ── Draw one lens at (cx, cy) ──
local function draw_lens(cx, cy, Rl, Rr, h, label)
    local t = auto_t(Rl, Rr, h)
    local absRl = math.abs(Rl)
    local absRr = math.abs(Rr)

    -- Left surface vertex at cx - t/2, right vertex at cx + t/2
    local vl = cx - t/2
    local vr = cx + t/2

    -- Left surface: circle of radius |Rl|
    -- If Rl > 0: center at vl + Rl (to the right), surface = left side of circle
    -- If Rl < 0: center at vl + Rl (to the left),  surface = right side of circle
    local cl_x = vl + Rl

    -- Right surface: circle of radius |Rr|
    -- If Rr > 0: center at vr - Rr (to the left),  surface = right side of circle -- wait
    -- Actually: right surface uses OPPOSITE convention since it faces the other way
    -- Rr < 0 means bulges RIGHT (convex outward) for the exit surface
    -- Rr > 0 means caves LEFT (concave) for exit surface
    -- Center at vr - Rr
    local cr_x = vr + Rr

    local N = 80
    local left_pts = {}
    local right_pts = {}

    for i = 0, N do
        local y = -h + 2*h*(i/N)
        -- Left surface x at height y
        local ql = math.sqrt(math.max(absRl*absRl - y*y, 0))
        local lx
        if Rl > 0 then
            lx = cl_x - ql   -- left side of circle (bulges left)
        else
            lx = cl_x + ql   -- right side of circle (caves right)
        end

        -- Right surface x at height y
        local qr = math.sqrt(math.max(absRr*absRr - y*y, 0))
        local rx
        if Rr < 0 then
            -- center is left of vertex, surface = right side of circle
            rx = cr_x + qr
        else
            -- center is right of vertex, surface = left side of circle
            rx = cr_x - qr
        end

        table.insert(left_pts,  {lx, cy + y})
        table.insert(right_pts, {rx, cy + y})
    end

    -- Fill
    local p = "\\fill[atlantic!12, opacity=0.85] "
    for i, v in ipairs(left_pts) do
        p = p .. (i==1 and "" or " -- ") .. string.format("(%.4f,%.4f)", v[1], v[2])
    end
    for i = #right_pts, 1, -1 do
        p = p .. string.format(" -- (%.4f,%.4f)", right_pts[i][1], right_pts[i][2])
    end
    tex.sprint(p .. " -- cycle;")

    -- Outlines
    for _, arc in ipairs({left_pts, right_pts}) do
        local s = "\\draw[atlantic!60, thin] "
        for i, v in ipairs(arc) do
            s = s .. (i==1 and "" or " -- ") .. string.format("(%.4f,%.4f)", v[1], v[2])
        end
        tex.sprint(s .. ";")
    end

    -- Top/bottom edges
    tex.sprint(string.format("\\draw[atlantic!60, thin] (%.4f,%.4f)--(%.4f,%.4f);",
        left_pts[1][1], left_pts[1][2], right_pts[1][1], right_pts[1][2]))
    tex.sprint(string.format("\\draw[atlantic!60, thin] (%.4f,%.4f)--(%.4f,%.4f);",
        left_pts[#left_pts][1], left_pts[#left_pts][2],
        right_pts[#right_pts][1], right_pts[#right_pts][2]))

    -- Optical axis tick
    tex.sprint(string.format("\\draw[warmgrey, dashed, very thin] (%.2f,%.2f)--(%.2f,%.2f);",
        cx - 1.5, cy, cx + 1.5, cy))

    -- Label below
    tex.sprint(string.format(
        "\\node[font=\\sffamily\\scriptsize, text=bk70, anchor=north] at (%.2f,%.2f) {%s};",
        cx, cy - h - 0.15, label))

    -- R values above
    local rl_str = (math.abs(Rl) > 500) and "\\infty" or string.format("%.0f", Rl)
    local rr_str = (math.abs(Rr) > 500) and "\\infty" or string.format("%.0f", Rr)
    tex.sprint(string.format(
        "\\node[font=\\sffamily\\tiny, text=warmgrey, anchor=south] at (%.2f,%.2f) {$R_L=%s$~~$R_R=%s$};",
        cx, cy + h + 0.1, rl_str, rr_str))
end

-- ════════════════════════════════════════════════════════════
-- GALLERY — 6 lens types in a 3x2 grid
-- ════════════════════════════════════════════════════════════
local col_spacing = 3.8
local row_spacing = 3.2

local lenses = {
    -- row 1
    {col=0, row=0, Rl= 3, Rr=-3,   label="Biconvex (symmetric)"},
    {col=1, row=0, Rl= 2, Rr=-6,   label="Biconvex (asymmetric)"},
    {col=2, row=0, Rl=-3, Rr= 3,   label="Biconcave (symmetric)"},
    -- row 2
    {col=0, row=1, Rl= 3, Rr= 3,   label="Meniscus (converging)"},
    {col=1, row=1, Rl= 999, Rr=-3,  label="Plano-convex"},
    {col=2, row=1, Rl= 999, Rr= 3,  label="Plano-concave"},
}

for _, L in ipairs(lenses) do
    local cx = L.col * col_spacing
    local cy = -L.row * row_spacing
    draw_lens(cx, cy, L.Rl, L.Rr, h, L.label)
end

\end{luacode}

\end{tikzpicture}
\end{document}
